% !TEX root = ../analisi-dei-capitolati.tex
% Capitolo 7 - Conclusioni

\section{Conclusioni}\label{sec:conclusioni}

La Tabella~\ref{tab:confronto} confronta i tre finalisti secondo i criteri
della Sezione~\ref{sec:metodologia}. \emph{Alto} è sempre il valore
favorevole: per i vincoli esterni indica vincoli contenuti.

\begin{table}[H]
  \centering
  \caption{Confronto tra i capitolati finalisti\label{tab:confronto}}
  \begin{tabularx}{\linewidth}{@{}Xccc@{}}
    \toprule
    \intestazione{Criterio} & \intestazione{C2 -- Concierge} &
    \intestazione{C1 -- AURA} & \intestazione{C0 -- AIRA} \\
    \midrule
    Interesse                & Alto  & Alto  & Medio \\
    Valore formativo         & Alto  & Alto  & Alto  \\
    Chiarezza del perimetro  & Alto  & Alto  & Medio \\
    Sostenibilità            & Medio & Medio & Basso \\
    Vincoli esterni          & Medio & Basso & Alto  \\
    Supporto del proponente  & Alto  & Alto  & Alto  \\
    \bottomrule
  \end{tabularx}
\end{table}

Il Git Push Team si candida per il capitolato \textbf{C2 -- NEXUM
Concierge} di Eggon S.r.l., che offre il miglior equilibrio tra
innovazione, perimetro chiaro e sostenibilità. In caso di mancata
aggiudicazione l'ordine di preferenza è \textbf{C1 -- AURA} e poi
\textbf{C0 -- AIRA}.
